%!TEX root = ../template.tex
%%%%%%%%%%%%%%%%%%%%%%%%%%%%%%%%%%%%%%%%%%%%%%%%%%%%%%%%%%%%%%%%%%%%
%% chapter4.tex
%% NOVA thesis document file
%%
%% Chapter with lots of dummy text
%%%%%%%%%%%%%%%%%%%%%%%%%%%%%%%%%%%%%%%%%%%%%%%%%%%%%%%%%%%%%%%%%%%%

\typeout{NT FILE chapter4.tex}%

\chapter{Conclusions}
\label{cha:conclusions}

In this report, it was possible to contextualize the scenario for the proposed thesis "Human Posture and Motion Classification through RF Signals", through a brief study of some modern technologies introduced in Chapter \ref{cha:introduction}.

In Chapter \ref{cha:related_work} it was possible to go into further detail about some of the technologies envisioned for the work to be pursued in this thesis. The fundamentals of FMCW operation were introduced, which brought to light the kind of technical gear that will be dealt with during the scope of this thesis. It was possible to understand how to gather range and velocity information, which is essential for classifying human motion.

From the multiple works already available in the literature, it was rapidly concluded that features played an important role in the behaviour of learning models, which led to a quick analysis of some of the possible candidates to implement for this thesis scenario, conducted in Chapter \ref{cha:related_work}. 
Given the vast range of machine-learning models available today, it was necessary to understand each category's work principle to select the most appropriate model for this thesis. 

Since there's some related work already conducted using the FMCW radar for classification of motion, the state of the art was also presented in Section \ref{sec:related_works}, where a summary of a selection of articles was introduced, which led to a better understanding of what future work could be developed.

Lastly, a proposed schedule was constructed in Chapter \ref{cha:work_plan}, segmenting the work of this thesis into six different subtasks, which will be the focus of the work to be developed in the next months. 
